\frfilename{mt123.tex}
\versiondate{18.11.04}
\copyrightdate{1994}

\def\chaptername{Integrasi}
\def\sectionname{Teorema-teorema konvergensi}

\newsection{123}

Jerih payah besar yang telah kita curahkan sejauh ini belum juga mendapat
pembenaran dari teorema mana pun yang cukup ampuh untuk membuatnya layak. Kini kita
sampai pada jantung teori integrasi modern, `teorema-teorema
konvergensi', yang menguraikan syarat-syarat yang memungkinkan kita
mengintegralkan limit suatu barisan fungsi terintegralkan.

\leader{123A}{Teorema B. Levi} Misalkan $(X,\Sigma,\mu)$ suatu ruang ukur
dan $\langle f_n\rangle_{n\in\Bbb N}$ suatu barisan fungsi bernilai
real, semuanya terintegralkan pada $X$, sedemikian sehingga (i)
$f_n\leae f_{n+1}$ untuk setiap $n\in\Bbb N$ (ii)
$\sup_{n\in\Bbb N}\int f_n<\infty$. Maka
$f=\lim_{n\to\infty}f_n$ terintegralkan, dan
$\int f=\lim_{n\to\infty}\int f_n$.

\medskip

\noindent{\bf Catatan}\cmmnt{ Perlu segera saya ulangi
kesepakatan-kesepakatan yang saya gunakan di sini. Setiap fungsi $f_n$
dianggap didefinisikan pada suatu himpunan koterabaikan
$\dom f_n\subseteq X$, seperti dalam 122Nc, dan fungsi limit $f$
dianggap mempunyai domain

\Centerline{$\{x:x\in\bigcup_{n\in\Bbb N}\bigcap_{m\ge n}\dom f_m$,
$\lim_{n\to\infty}f_n(x)$ ada di $\Bbb R\}$,}

\noindent seperti dalam 121Fa.} \dvro{Pernyataan}{Anda tidak akan
kehilangan gagasan penting apa pun jika mengandaikan bahwa setiap
$f_n$ didefinisikan di seluruh $X$; tetapi pernyataan}
`$f$ terintegralkan' mencakup pernyataan `$f$ didefinisikan, sebagai
bilangan real, hampir di mana-mana'\cmmnt{, dan ini merupakan bagian
hakiki dari teorema tersebut}.

\proof{{\bf (a)} Mula-mula kita tangani kasus ketika
$f_0=0$ a.e.
Tetapkan $c=\sup_{n\in\Bbb N}\int f_n=\lim_{n\to\infty}\int f_n$
(dengan memperhatikan bahwa, menurut 122Od,
$\sequencen{\int f_n}$ merupakan barisan tak-menurun).

\medskip

\quad{\bf (i)} Semua himpunan $\dom f_n$, $\{x:f_0(x)=0\}$,
$\{x:f_n(x)
\le f_{n+1}(x)\}$ bersifat koterabaikan, sehingga irisannya, $F$, juga
koterabaikan. Untuk setiap $n\in\Bbb N$ terdapat himpunan
koterabaikan $E_n$ sedemikian sehingga $f_n\restr E_n$ terukur (122P);
misalkan $E^*$ suatu himpunan terukur koterabaikan yang termuat dalam
himpunan koterabaikan $F\cap\bigcap_{n\in\Bbb N}E_n$.

\medskip

\quad{\bf (ii)} Untuk $a>0$ dan $n\in\Bbb N$, tetapkan
$H_n(a)=\{x:x\in E^*$, $f_n(x)\ge a\}$;
maka $H_n(a)$ terukur karena $f_n\restr E_n$ terukur dan $E^*$
merupakan subhimpunan terukur dari $E_n$. Selain itu, $a\chi
H_n(a)
\le f_n$ di setiap titik pada $E^*$, sehingga

\Centerline{$a\mu H_n(a)=\int a\chi H_n(a)\le\int f_n\le c$.}

\noindent Karena $f_n(x)\le f_{n+1}(x)$ untuk setiap $x\in E^*$,
$H_n(a)\subseteq H_{n+1}(a)$ untuk setiap $n\in\Bbb N$, dan dengan
menuliskan $H(a)=\bigcup_{n\in\Bbb N}H_n(a)$, kita mempunyai

\Centerline{$\mu H(a)=\lim_{n\to\infty}\mu H_n(a)\le\Bover{c}{a}$}

\noindent (112Ce). Khususnya, $\mu H(a)<\infty$ untuk setiap $a$.
Selain itu,

\Centerline{$\mu(\bigcap_{k\ge 1}H(k))\le\inf_{k\ge 1}\mu H(k)
\le\inf_{k\ge 1}\Bover{c}{k}=0$.}

\noindent Tetapkan $E=E^*\setminus\bigcap_{k\ge 1}H(k)$; maka $E$
bersifat koterabaikan.

\medskip

\quad{\bf (iii)} Jika $x\in E$, terdapat suatu $k$ sedemikian sehingga
$x\notin H(k)$, yakni $x\notin\bigcup_{n\in\Bbb N}H_n(k)$, yakni
$f_n(x)<k$ untuk setiap $n$; selain itu,
$\langle f_n(x)\rangle_{n\in\Bbb N}$ merupakan barisan tak-menurun,
sehingga
$f(x)=\lim_{n\to\infty}f_n(x)=\sup_{n\in\Bbb N}f_n(x)$ ada di
$\Bbb R$. Dengan demikian, fungsi
limit $f$ didefinisikan hampir di mana-mana.
Karena setiap $f_n\restr E$ terukur (121Eh), demikian pula
$f\restr E=\lim_{n\to\infty}f_n\restr E$ (121Fa).
Jika $\epsilon>0$, maka $\{x:x\in E,\,f(x)\ge\epsilon\}$ termuat dalam
$H({1\over 2}\epsilon)$, sehingga mempunyai ukuran berhingga.

\medskip

\quad{\bf (iv)}
Sekarang andaikan bahwa $g$ suatu fungsi sederhana dan $g\leae f$.
Seperti dalam bukti 122G, $H=\{x:g(x)\ne 0\}$ mempunyai ukuran
berhingga, dan $g$ dibatasi di atas oleh suatu $M$.

Misalkan $\epsilon>0$. Untuk setiap $n\in\Bbb N$, tetapkan
$G_n=\{x:x\in E,\,(g-f_n) (x)\ge\epsilon\}$.
Maka setiap $G_n$ terukur, dan $\langle G_n \rangle_{n\in\Bbb N}$
merupakan barisan tak-menaik dengan irisan

\Centerline{$\{x:x\in E,\,g(x)\ge\epsilon +\sup_{n\in\Bbb N}f_n(x)\}
\subseteq\{x:g(x)>f(x)\}$,}

\noindent yang terabaikan.
Selain itu, $\mu G_0<\infty$ karena $G_0\subseteq H$.
Akibatnya, $\lim_{n\to\infty}\mu G_n=0$ (112Cf).
Ambil $n$ sedemikian sehingga $\mu G_n\le\epsilon$. Maka, untuk setiap
$x\in E$,

\Centerline{$g(x)\le f_n(x)+\epsilon\chi H(x)+M\chi G_n(x)$,}

\noindent sehingga

\Centerline{$g\leae f_n+M\chi G_n+\epsilon\chi H$}

\noindent dan

\Centerline{$\int g\le\int f_n+M\mu G_n+\epsilon\mu H
\le c+\epsilon(M+\mu H)$.\footnote{Saya berterima kasih kepada
P. Wallace Thompson karena telah menemukan kekeliruan pada tahap ini
dalam edisi-edisi sebelumnya.}}

\noindent Karena $\epsilon$ sebarang, $\int g\le c$.

\medskip

\quad{\bf (v)} Dengan demikian, $f\restr E$ (yang nonnegatif)
memenuhi syarat-syarat Lema 122Ja, dan terintegralkan. Selain itu,
integralnya paling besar $c$, menurut Definisi 122K. Karena
$f\eae f\restr E$, $f$ juga terintegralkan, dengan integral yang sama
(122Rb). Di pihak lain, $f\geae f_n$ untuk setiap $n$, sehingga
$\int f\ge\sup_{n\in\Bbb N}\int f_n=c$, menurut 122Od.

Dengan ini kasus $f_0=0$ a.e.\ selesai dibuktikan.

\medskip

{\bf (b)} Untuk kasus umum, tinjau barisan
$\langle f'_n\rangle_{n\in\Bbb N}=\langle f_n-f_0\rangle_{n\in\Bbb N}$.
Menurut (a), $f'=\lim_{n\to\infty}f'_n$ terintegralkan, dan
$\int f'=\lim_{n\to\infty}
\int f'_n$; sekarang
$\lim_{n\to\infty}f_n\eae f'+f_0$, sehingga fungsi itu terintegralkan,
dengan integral
$\int f'+\int f_0=\lim_{n\to\infty}\int f_n$.

\medskip

\noindent{\bf Catatan} Mungkin Anda telah memperhatikan, tanpa merasa
heran, bahwa argumen (a-iv) dalam bukti ini mengulangi argumen yang
digunakan untuk kasus khusus 122G.
}%end of proof of 123A

\leader{123B}{Lema Fatou} Misalkan $(X,\Sigma,\mu)$ suatu ruang ukur,
dan $\langle f_n\rangle_{n\in\Bbb N}$ suatu barisan fungsi bernilai
real, semuanya terintegralkan pada $X$. Andaikan setiap $f_n$
nonnegatif a.e., dan
$\liminf_{n\to\infty}\int f_n<\infty$. Maka
$\liminf_{n\to\infty}
f_n$ terintegralkan, dan
$\int\liminf_{n\to\infty}f_n\le\liminf_{n\to\infty}\int f_n$.

\cmmnt{\medskip

\noindent{\bf Catatan} Sekali lagi, teorema ini mencakup pernyataan
bahwa $\liminf_{n\to\infty}f_n(x)$ ada di $\Bbb R$ untuk hampir setiap
$x\in X$.}

\proof{ Tetapkan $c=\liminf_{n\to\infty}\int f_n$ dan
$f=\liminf_{n\to\infty}f_n$.
Untuk setiap $n\in\Bbb N$, misalkan $E_n$ suatu himpunan koterabaikan
sedemikian sehingga $f_n'=f_n\restr E_n$ terukur dan nonnegatif.
Tetapkan $g_n=\inf_{m\ge n}f'_m$; maka setiap $g_n$ terukur (121Fc),
nonnegatif, dan didefinisikan pada himpunan koterabaikan
$\bigcap_{m\ge n}E_m$, serta $g_n\leae f_n$, sehingga $g_n$
terintegralkan (122P) dengan
$\int g_n\le\inf_{m\ge n}\int f_m\le c$.
Selanjutnya, $g_n(x)\le g_{n+1}(x)$ untuk setiap
$x\in\dom g_n$, sehingga $\langle g_n\rangle_{n\in\Bbb N}$ memenuhi
syarat-syarat teorema B. Levi (123A), dan
$g=\lim_{n\to\infty}g_n$ terintegralkan, dengan
$\int g=\lim_{n\to\infty}\int g_n\le c$.
Akhirnya, karena setiap $f'_n$ sama dengan $f_n$ hampir di mana-mana,
$g=\liminf_{n\to\infty}f'_n\eae f$, dan $\int f$ ada, sama dengan
$\int g\le c$.
}%end of proof of 123B

\leader{123C}{Teorema Konvergensi Terdominasi Lebesgue} Misalkan
$(X,\Sigma,\mu)$ suatu ruang ukur dan
$\langle f_n\rangle_{n\in\Bbb N}$ suatu barisan fungsi bernilai real,
semuanya terintegralkan pada $X$, sedemikian sehingga
$f(x)=\lim_{n\to\infty}f_n(x)$ ada di $\Bbb R$ untuk hampir setiap
$x\in X$. Andaikan pula bahwa terdapat suatu fungsi terintegralkan $g$
sedemikian sehingga $|f_n|\leae g$ untuk setiap $n$.
Maka $f$ terintegralkan, dan $\lim_{n\to\infty}\int f_n$ ada serta
sama dengan $\int f$.

\proof{ Tinjau $\tilde f_n=f_n+g$ untuk setiap $n\in\Bbb N$.
Maka $0\le\tilde f_n\le 2g$ a.e.\ untuk setiap $n$, sehingga
$\tilde c=\liminf_{n\to\infty}\int \tilde f_n$
ada di $\Bbb R$, dan $\tilde f=\liminf_{n\to\infty}\tilde f_n$
terintegralkan, dengan $\int \tilde f\le \tilde c$, menurut Lema Fatou
(123B).
Namun, perhatikan bahwa $f\eae\tilde f-g$, karena
$f(x)=\tilde f(x)-g(x)$ setidaknya setiap kali $f(x)$ dan $g(x)$
keduanya terdefinisi, sehingga $f$ terintegralkan, dengan

\Centerline{$\int f=\int \tilde f-\int g
\le\liminf_{n\to\infty}\int \tilde f_n-\int g
=\liminf_{n\to\infty}\int f_n$.}

\noindent Demikian pula, dengan meninjau
$\langle -f_n\rangle_{n\in\Bbb N}$,

\Centerline{$\int(-f)\le\liminf_{n\to\infty}\int(-f_n)$,}

\noindent yakni,

\Centerline{$\int f\ge\limsup_{n\to\infty}\int f_n$.}

\noindent Jadi $\lim_{n\to\infty}\int f_n$ ada dan sama dengan
$\int f$.
}%end of proof of 123C

\cmmnt{\medskip

\noindent{\bf Catatan} Akhirnya kita telah sampai pada tahap ketika
masalah-masalah teknis yang berkaitan dengan fungsi-fungsi terdefinisi
sebagian mulai berkurang, atau lebih tepatnya, tertangani secara
efisien oleh kesepakatan-kesepakatan yang saya gunakan mengenai
penafsiran rumus-rumus seperti `$\limsup$'.
}%end of comment

\leader{123D}{}\cmmnt{ Sebagai gambaran tentang kekuatan
teorema-teorema ini, saya berikan di sini suatu hasil yang merupakan
salah satu penerapan standar teori tersebut.

\medskip

\noindent}{\bf Korolari} Misalkan $(X,\Sigma,\mu)$ suatu ruang ukur dan
$\ooint{a,b}$ suatu interval terbuka tak kosong di dalam $\Bbb R$.
Misalkan $f:X\times\ooint{a,b}\to\Bbb R$ suatu fungsi sedemikian
sehingga

\inset{(i) integral $F(t)=\int f(x,t)dx$ terdefinisi untuk setiap
$t\in\ooint{a,b}$;}

\inset{(ii) turunan parsial $\pd{f}{t}$ dari $f$ terhadap variabel
kedua terdefinisi di setiap titik dalam $X\times\ooint{a,b}$;}

\inset{(iii) terdapat fungsi terintegralkan
$g:X\to\coint{0,\infty}$ sedemikian sehingga
$|\pd{f}{t}(x,t)|\le g(x)$ untuk setiap $x\in X$ dan
$t\in\ooint{a,b}$.}

\noindent Maka turunan $F'(t)$ dan integral
$\biggerint\pd{f}{t}(x,t)dx$ ada untuk setiap $t\in\ooint{a,b}$,
dan keduanya sama.

\proof{{\bf (a)} Misalkan $t$ sembarang titik dalam $\ooint{a,b}$, dan
$\sequencen{t_n}$ sembarang barisan dalam
$\ooint{a,b}\setminus\{t\}$ yang konvergen ke $t$. Tinjau

\Centerline{$\Bover{F(t_n)-F(t)}{t_n-t}
=\int\Bover{f(x,t_n)-f(x,t)}{t_n-t}\mu(dx)$}

\noindent untuk setiap $n$. (Langkah ini menggunakan 122O.) Jika kita
menetapkan

\Centerline{$f_n(x)=\Bover{f(x,t_n)-f(x,t)}{t_n-t}$,}

\noindent untuk $x\in X$, maka dari Teorema Nilai Rata-Rata kita
melihat bahwa terdapat suatu $\tau$ (yang tentu saja bergantung pada
$n$ maupun $x$), yang terletak di antara $t_n$ dan $t$, sedemikian
sehingga $f_n(x)=\pd{f}{t}(x,\tau)$, sehingga $|f_n(x)|\le g(x)$.
Pada saat yang sama,
$\lim_{n\to\infty}f_n(x)=\pd{f}{t}(x,t)$ untuk setiap $x$.
Jadi Teorema Konvergensi Terdominasi Lebesgue menunjukkan bahwa
$\biggerint\pd{f}{t}(x,t)dx$ ada dan sama dengan

\Centerline{$\lim_{n\to\infty}\int f_n(x)dx
=\lim_{n\to\infty}\Bover{F(t_n)-F(t)}{t_n-t}$.}

\wheader{123D}{6}{2}{2}{12pt}

{\bf (b)} Karena $\sequencen{t_n}$ sebarang,

\Centerline{$\lim_{s\to t}\Bover{F(s)-F(t)}{s-t}
=\int\Pd{f}{t}(x,t)dx$,}

\noindent sebagaimana dinyatakan.
}%end of proof of 123D

\cmmnt{\medskip

\noindent{\bf Catatan} Dalam jilid berikutnya saya menawarkan suatu
variasi teorema ini, dengan hipotesis maupun kesimpulannya diperlemah
(252Ye).
}%end of comment

\exercises{
\leader{123X}{Latihan dasar $\pmb{>}$(a)}
%\spheader 123Xa
Misalkan $(X,\Sigma,\mu)$ suatu ruang ukur, dan $\sequencen{f_n}$ suatu
barisan fungsi bernilai real, semuanya terintegralkan pada $X$,
sedemikian sehingga $\sum_{n=0}^{\infty}\int|f_n|$ berhingga.
Tunjukkan bahwa $f(x)=\sum_{n=0}^{\infty}f_n(x)$ ada di $\Bbb R$
untuk hampir setiap $x\in X$, dan bahwa
$\int f=\sum_{n=0}^{\infty}\int f_n$.
\Hint{tinjau terlebih dahulu kasus ketika setiap $f_n$ nonnegatif.}

\spheader 123Xb Misalkan $(X,\Sigma,\mu)$ suatu ruang ukur.
Andaikan $T$ sembarang subhimpunan $\Bbb R$, dan
$\langle f_t\rangle_{t\in T}$ suatu keluarga fungsi, semuanya
terintegralkan pada $X$, sedemikian sehingga, untuk setiap $t\in T$,

\Centerline{$f_t(x)=\lim_{s\in T,s\to t}f_s(x)$}

\noindent untuk hampir setiap $x\in X$. Andaikan pula bahwa terdapat
suatu fungsi terintegralkan $g$ sedemikian sehingga $|f_t|\leae g$
untuk setiap $t\in T$.
Tunjukkan bahwa $t\mapsto\int f_t:T\to\Bbb R$ kontinu.

\sqheader 123Xc Misalkan $f$ suatu fungsi bernilai real yang
didefinisikan di setiap titik pada $\coint{0,\infty}$, dan lengkapi
interval ini dengan ukuran Lebesgue. {\bf Transformasi Laplace} (real)
darinya ialah fungsi $F$ yang didefinisikan oleh

\Centerline{$F(s)=\int_0^{\infty}e^{-sx}f(x)dx$}

\noindent untuk semua bilangan real $s$ yang membuat integral tersebut
terdefinisi.

\quad(i) Tunjukkan bahwa jika $s\in\dom F$ dan $s'\ge s$, maka
$s'\in\dom F$ (karena $e^{-s'x}e^{sx}\le 1$ untuk setiap $x$).
(Bagaimana Anda mengetahui bahwa $x\mapsto e^{-s'x}e^{sx}$ terukur?)

\quad(ii) Tunjukkan bahwa $F$ dapat didiferensialkan pada bagian dalam
domainnya.
\Hint{perhatikan bahwa jika $a_0\in\dom F$ dan $a_0<a<b$, maka
terdapat suatu $M$ sedemikian sehingga
$xe^{-sx}|f(x)|\le Me^{-a_0x}|f(x)|$ setiap kali
$x\in\coint{0,\infty}$, $s\in[a,b]$.}

\quad(iii) Tunjukkan bahwa jika $F$ terdefinisi untuk setidaknya satu
nilai, maka
$\lim_{s\to\infty}F(s)=0$. \Hint{gunakan Teorema Konvergensi
Terdominasi Lebesgue untuk menunjukkan bahwa
$\lim_{n\to\infty}F(s_n)=0$ setiap kali
$\lim_{n\to\infty}s_n=\infty$.}

\quad(iv) Tunjukkan bahwa jika $f$, $g$ mempunyai transformasi Laplace
$F$, $G$, maka transformasi Laplace dari $f+g$ adalah $F+G$,
setidaknya pada $\dom F\cap \dom
G$.

\spheader 123Xd Misalkan $(X,\Sigma,\mu)$ suatu ruang ukur
dan $\langle f_n\rangle_{n\in\Bbb N}$ suatu barisan fungsi bernilai
real, semuanya terintegralkan pada $X$, sedemikian sehingga terdapat
suatu fungsi terintegralkan $g$ dengan $|f_n|\leae g$ untuk setiap $n$.
Tunjukkan bahwa $\limsup_{n\to\infty}f_n$ terintegralkan dan bahwa
$\int\limsup_{n\to\infty}f_n\ge\limsup_{n\to\infty}\int f_n$.
%123C

\leader{123Y}{Latihan lanjutan (a)}
%\spheader 123Ya also in \S235
Misalkan $(X,\Sigma,\mu)$ suatu ruang ukur, $Y$ sembarang himpunan dan
$\phi:X\to Y$ sembarang fungsi; misalkan $\mu\phi^{-1}$ ukuran citra
pada $Y$ (112Xf). Tunjukkan bahwa jika $h:Y\to\Bbb R$ sembarang
fungsi, maka $h$ terintegralkan terhadap $\mu\phi^{-1}$ jika dan hanya
jika $h\phi$ terintegralkan terhadap $\mu$; dalam hal itu kedua integral
tersebut sama.

\header{123Yb}{\bf (b)} Jelaskan cara menyesuaikan 123Xc dengan kasus
ketika $f$ tidak didefinisikan pada suatu subhimpunan terabaikan dari
$\Bbb R$.

\header{123Yc}{\bf (c)} Misalkan $(X,\Sigma,\mu)$ suatu ruang ukur dan
$a<b$ dalam $\Bbb R$. Misalkan
$f:X\times\ooint{a,b}\to\coint{0,\infty}$ suatu fungsi sedemikian
sehingga $\int f(x,t)dx$ terdefinisi untuk setiap $t\in\ooint{a,b}$
dan $t\mapsto f(x,t)$ kontinu untuk setiap $x\in X$. Andaikan bahwa
$c\in\ooint{a,b}$ sedemikian sehingga
$\liminf_{t\to c}\int f(x,t)dx<\infty$.
Tunjukkan bahwa $\int\liminf_{t\to c}f(x,t)dx$ terdefinisi dan kurang
dari atau sama dengan $\liminf_{t\to c}\int f(x,t)dx$.

\header{123Yd}{\bf (d)} Tunjukkan bahwa terdapat suatu fungsi
$f:\BbbR^2\to\{0,1\}$ sedemikian sehingga (i) integral Lebesgue
$\int f(x,t)dx$ terdefinisi dan sama dengan $1$ untuk setiap $t\ne 0$
(ii) fungsi $x\mapsto\liminf_{t\to 0}f(x,t)$ tidak terukur Lebesgue.
({\it Catatan\/}: tentu saja Anda harus memulai konstruksi dari suatu
subhimpunan tak-terukur dari $\Bbb R$; lihat 134B untuk himpunan
semacam itu.)

\spheader 123Ye Misalkan $(Y,\Tau,\nu)$ suatu ruang ukur. Misalkan $X$
suatu himpunan, $\Sigma$ suatu aljabar-$\sigma$ atas
subhimpunan-subhimpunan $X$, dan
$\langle\mu_y\rangle_{y\in Y}$ suatu keluarga ukuran pada $X$
sedemikian sehingga $\mu_yX$ berhingga untuk setiap $y$ dan
$\mu E=\int\mu_yE\,\nu(dy)$ terdefinisi untuk setiap $E\in\Sigma$.
(i) Tunjukkan bahwa $\mu:\Sigma\to\coint{0,\infty}$ merupakan suatu
ukuran. (ii) Tunjukkan bahwa jika
$f:X\to\coint{0,\infty}$ suatu fungsi $\Sigma$-terukur, maka $f$
terintegralkan terhadap $\mu$ jika dan hanya jika fungsi itu
terintegralkan terhadap $\mu_y$ untuk hampir setiap $y\in Y$ dan
$\int\bigl(\int fd\mu_y\bigr)\nu(dy)$ terdefinisi, serta bahwa nilai
ini dalam hal itu sama dengan $\int fd\mu$.

\spheader 123Yf Misalkan $(X,\Sigma,\mu)$ suatu ruang ukur, dan
$\sequencen{f_n}$ suatu barisan fungsi bernilai real yang terukur
secara virtual dan semuanya didefinisikan hampir di mana-mana dalam
$X$. Andaikan bahwa
$\sum_{n=0}^{\infty}\int|f_n(x)-1|\mu(dx)<\infty$.
Tunjukkan bahwa $\prod_{n=0}^{\infty}f_n(x)$ ada di $\Bbb R$ untuk
hampir setiap $x\in X$.
%123A

}%end of exercises

\endnotes{
\Notesheader{123} Saya berharap 123D dan kasus khususnya, 123Xc, akan
membantu meyakinkan Anda bahwa teori di sini mempunyai penerapan yang
berguna.

Semua teorema dalam bagian ini dapat dipandang sebagai teorema
`pertukaran limit', yang menetapkan syarat-syarat yang memungkinkan

$$\lim_{n\to\infty}\int f_n=\int\lim_{n\to\infty}f_n,$$

\noindent atau

$$\Bover{\partial}{\partial t}{\int}f\,dx=\int\Bover{\partial
f}{\partial
t}dx.$$

\noindent Bahkan untuk fungsi-fungsi yang dapat ditangani dengan metode
integrasi yang jauh lebih elementer (misalnya, integral Riemann),
teorema-teorema jenis ini dapat menuntut verifikasi
pertidaksamaan-pertidaksamaan yang melelahkan. Kekuatan integral
Lebesgue terletak pada teorema-teorema umum yang diberikannya, yang
mencakup bagian yang cukup besar dari kasus-kasus penting yang
muncul dalam praktik. (Namun harus saya akui bahwa tidak ada yang lebih
khas dalam analisis terapan daripada kebutuhannya akan hasil-hasil
khusus yang berkaitan dengan, tetapi tidak dapat diturunkan dari,
teorema-teorema umum yang standar.) Sebagai contoh, dalam 123Xc, fakta
bahwa daerah integrasinya merupakan interval tak terbatas
$\coint{0,\infty}$ tidak menambah kesulitan. Tentu saja, ini berkaitan
dengan fakta bahwa kita hanya meninjau integral fungsi-fungsi yang
nilai mutlaknya terintegralkan.

Limit-limit yang digunakan dalam 123A-123C semuanya merupakan limit
barisan; tentu saja, bagian dari hakikat teori ukur ialah bahwa kita
berharap dapat menangani keluarga terhitung dari himpunan atau fungsi,
tetapi keluarga yang lebih besar dari itu menimbulkan kekhawatiran. Meskipun
demikian, terdapat banyak konteks yang memungkinkan kita mengambil
jenis limit lain. Saya menguraikan beberapa di antaranya dalam 123D,
123Xb, dan 123Xc(iii). Intinya ialah bahwa dalam limit seperti
$\lim_{t\to u}\phi(t)$, dengan $u\in[-\infty,\infty]$, kita akan
mempunyai $\lim_{t\to u}\phi(t)=a$ jika dan hanya jika
$\lim_{n\to\infty}\phi(t_n)=a$ setiap kali $\sequencen{t_n}$ konvergen
ke $u$; sehingga ketika mencari suatu limit
$\lim_{t\to u}\int f_t$, untuk suatu keluarga fungsi
$\langle f_t\rangle_{t\in T}$, cukuplah jika kita dapat menentukan
$\lim_{n\to\infty}\int f_{t_n}$ untuk cukup banyak barisan
$\sequencen{t_n}$. Argumen jenis ini akan efektif bagi setiap limit
standar $\lim_{t\uparrow a}$, $\lim_{t\downarrow a}$,
$\lim_{t\to a}$, $\lim_{t\to\infty}$, $\lim_{t\to-\infty}$ dari
kalkulus dasar, dan dapat digunakan bersama dengan teorema B. Levi
ataupun teorema Lebesgue. Barangkali perlu saya catat bahwa suatu
kesulitan muncul dalam perluasan serupa dari lema Fatou
(123Yc-123Yd).
}%end of notes

\frnewpage
